\documentclass{article}
\usepackage[UTF8]{ctex}
\usepackage{amsmath,amssymb,amsfonts}
\usepackage{enumitem}
\usepackage{geometry}
\geometry{a4paper, margin=1in}
% === 错题本排版优化（不参与题目标识） ===
\setlength{\parindent}{0pt}
\setlength{\parskip}{0.35em plus 0.10em minus 0.08em}
\setlength{\emergencystretch}{3em}
\setlist{itemsep=0.25em, topsep=0.40em, parsep=0pt}
\allowdisplaybreaks[2]
% === 错题本排版优化结束 ===
\begin{document}
\section*{东华大学2026年数学分析真题}
\section*{49．东华大学2026年数学分析试卷}

\subsection{题目1}
1．求极限 \(\lim _{n \rightarrow \infty} \sum_{i=1}^{n} \sin \left(\frac{i}{n^{2}}\right) \ln \left(1+\frac{i}{n}\right)\) ．

\subsection{题目2}
2．求极限 \(\lim _{x \rightarrow 0} \frac{1}{x}\left(\frac{1}{x}-\cot x\right)\) 。

\subsection{题目3}
3 ．求积分 \(\int_{0}^{1} \frac{\ln x}{1-x} \mathrm{~d} x\) ．

\subsection{题目4}
4．求幂级数 \(\sum_{n=0}^{\infty} \frac{x^{2 n+2}}{(n+1)(2 n+1)}\) 的收敛域及和函数。

\subsection{题目5}
5．求积分 \(\oint_{\sum} \frac{x \mathrm{~d} y \mathrm{~d} z+y \mathrm{~d} z \mathrm{~d} x+z \mathrm{~d} x \mathrm{~d} y}{\left(x^{2}+y^{2}+z^{2}\right)^{\frac{3}{2}}}\) ，其中 \(\sum\) 为椭球面 \(x^{2}+2 y^{2}+3 z^{2}=6\) 的内侧．

二．（15 分）设实数列 \(\left\{a_{n}\right\}\) 满足 \(a_{1}>0, a_{n+1}=a_{n}+\frac{1}{n a_{n}}, n=1,2, \cdots\) ．证明： \(\lim _{n \rightarrow \infty} \frac{a_{n}^{2}}{\ln n}=2\) ．

三．（20 分）叙述实数完备性的有限覆盖定理，并用之证明一元函数的介值性定理．

四．（15 分）设函数 \(f(x)\) 在 \((-\infty,+\infty)\) 上具有二阶连续导数，证明：\(f^{\prime \prime}(x) \geq 0\) 的充要条件是对任意不同实数 \(a, b\) ，有 \(f\left(\frac{a+b}{2}\right) \leq \frac{1}{b-a} \int_{a}^{b} f(x) \mathrm{d} x\) ．

\subsection{题目6}

\subsection{题目1}
原题第1题：讨论反常积分 \(\int_{0}^{+\infty} \frac{\sin \sqrt{x}}{x^{p}} \mathrm{~d} x\) 的敛散性，其中参数 \(p>\frac{1}{2}\) ．

\subsection{题目7}

\subsection{题目2}
原题第2题：讨论函数列 \(f_{n}(x)=\frac{n x}{n x+1}, n=1,2, \cdots, x \in[0,+\infty)\) 的一致收敛性及极限函数的连续性和可微性．

\subsection{题目8}

\subsection{题目3}
原题第3题：讨论级数 \(\sum_{n=0}^{\infty}\left|\sin \left(\pi \sqrt{n^{2}+1}\right)\right|^{p}\) 的敛散性，其中 \(p>0\) ．

六．（15 分）取 \(u=\frac{x}{x^{2}+y^{2}}, v=\frac{y}{x^{2}+y^{2}}\) 为新的自变量，变换方程 \(\frac{\partial^{2} z}{\partial x^{2}}+\frac{\partial^{2} z}{\partial y^{2}}=0\) ，其中所涉及的各阶导数均存在且连续。
\end{document}
